\documentclass[../../main.tex]{subfiles}
\begin{document}

\begin{center}
  {\fontsize{48pt}{58pt}\selectfont \textbf{2011}\normalsize\textbf{年度}}
\end{center}

\begin{table}[htb]
  \centering
  \caption{出題テーマと難易度}
  \begin{tabular}{|c|p{9cm}|c|}
    \hline
    \textbf{問題番号} & \textbf{テーマ} & \textbf{難易度} \\
    \hline
    第1問 & 行列による一次変換，不変直線と面積・級数 & \\
    \hline
    第2問 & 絶対値を含む定積分の最小値と積分 & \\
    \hline
    第3問 & 円と直線上の点がつくる三角形の面積の最大値 & \\
    \hline
    第4問 & 正方形を回転して得られる立体の体積の最大値 & \\
    \hline
  \end{tabular}
\end{table}

\begin{table}[htb]
  \centering
  \caption{解答}
  \renewcommand{\arraystretch}{1.4}
  \begin{tabular}{|c|c|p{9cm}|}
    \hline
    \textbf{問題番号} & \textbf{小問} & \textbf{答え} \\
    \hline
    第1問 & (1)〜(3) & \\
    \hline
    第2問 & (1) & 最小値 $\sqrt{2}-1$ \\
    \cline{2-3}
    & (2) & $\displaystyle\int_0^1 f(x)\,dx=\dfrac{1}{4}+\dfrac{1}{2}\log 2$ \\
    \hline
    第3問 & -- & 最大値 $\dfrac{k-1}{2(k+1)}$ \\
    \hline
    第4問 & (1) & 軸は正方形とただ1点で交わる（証明問題） \\
    \cline{2-3}
    & (2) & 最大値 $\sqrt{2}\pi$ \\
    \hline
  \end{tabular}
\end{table}

\paragraph{試験形態}
4問．試験時間は要確認．

\end{document}
